\begin{definition}[The correlated charge state]
    \label{def:peps_regular_correlated_charge_state}
    \lean{TNLean.PEPS.regularCorrelatedChargeState}
    \leanok
    \uses{def:peps_regular_edge_charge_contraction}
    Let $e_0,e_1$ be ordered edges of a finite graph, and let $u$ be a
    bond assignment. Resolve the remote label $q$ by inserting
    $\mathbf1_{\eta_{e_1}=q}$ on $e_1$, and put
    $\chi(pq^{-1}z\eta_{e_0})$ on $e_0$. Sum the resulting
    closed tensor states over $q$. Denote this correlated state by
    $\Psi_z(u;\chi,p)$. This definition retains the same original
    site coefficients in every summand. It is auxiliary to the literal
    charge-pair formula in \cite{Schuch2010PEPS}, lines 2505--2535.
\end{definition}

\begin{theorem}[Correlated weights in the physical cut]
    \label{thm:peps_regular_correlated_charge_cut}
    \lean{TNLean.PEPS.openRegionWeight_regularEdgeCharacterSite_eq_weighted,
        TNLean.PEPS.regularCorrelatedChargeState_apply,
        TNLean.PEPS.regionLocalTerm_mulVec_regularCorrelatedChargeState}
    \leanok
    \uses{def:peps_regular_correlated_charge_state,
        def:peps_regular_weighted_open_contraction,
        thm:peps_regular_coherent_global_transport}
    The preceding closed state has coefficient
    \[
        \Psi_z(u;\chi,p)(\sigma)
        =\sum_{\eta:E\to G}\chi(p\eta_{e_1}^{-1}z\eta_{e_0})
        \prod_v a_v\bigl(\eta^{u}|_v,\sigma_v\bigr).
    \]
    A diagonal inserted at an edge tail in a region gives precisely the
    corresponding weighted open coefficient. Suppose $t(e_0)\in R$
    and $t(e_1)\notin R$. If a regional matrix changes every actual
    weighted boundary column by
    \[
        \chi(pq^{-1}\eta_{e_0})
        \longmapsto\chi(pq^{-1}k^{-1}\eta_{e_0}),
    \]
    uniformly in $\chi,p,q$ and the crossing labels, its extension by
    the identity gives
    \[
        (W_R\otimes I)\Psi_1(u;\chi,p)=\Psi_{k^{-1}}(u;\chi,p).
    \]
    This is a linear contraction identity for
    \cite{Schuch2010PEPS}, lines 2505--2535 and 2560--2581.
    The column premise in this auxiliary statement is discharged by the
    actual native operation below.
\end{theorem}
\begin{proof}\leanok
    The product of the diagonal factors reduces to the factor at its
    actual tail vertex. Summing the remote indicator fixes
    $q=\eta_{e_1}$ and restores the joint character weight.
    For each fixed $q$, use the native physical cut identity; the remote
    diagonal belongs to the common complementary tensor. Pass the
    column action through that cut sum and then sum over $q$.
\end{proof}

\begin{definition}[The native remote charge edge]
    \label{def:peps_torus_charge_flux_partner_edge}
    \lean{TNLean.PEPS.torusChargeFluxPartnerEdge}
    \leanok
    \uses{def:peps_torus_charge_flux_crossing_strip}
    At $v=(x,y)$, the remote reference is the native horizontal bond
    joining $(x+1,y)$ to $(x+2,y)$, below the eight-site crossing strip.
    Both periodic seam orientations are retained.
    This is auxiliary native geometry for
    \cite{Schuch2010PEPS}, lines 2505--2581.
\end{definition}

\begin{theorem}[A physical crossing with the actual remote partner]
    \label{thm:peps_torus_correlated_charge_flux_crossing_global}
    \lean{TNLean.PEPS.torusChargeFluxPartnerEdge_tail_not_mem,
        TNLean.PEPS.IsGIsometric.exists_unitary_torusCorrelatedChargeFluxCrossing_global}
    \leanok
    \uses{def:peps_torus_charge_flux_partner_edge,
        thm:peps_regular_correlated_charge_cut,
        thm:peps_torus_charge_flux_crossing}
    Let $G$ be finite, let the periods satisfy $w\geq5$ and $H\geq4$,
    and let the native four-leg tensor be regular $G$-isometric.
    For every translated crossing strip there is a regional unitary
    $W_R$, chosen before $k,\chi,p$, whose extension by the complementary
    identity is unitary and satisfies
    \[
        (W_R\otimes I)\Psi_1(u^{(0)}_k;\chi,p)
        =\Psi_{k^{-1}}(u^{(0)}_k;\chi,p).
    \]
    Here $e_0$ is the displayed charge bond, $e_1$ is the actual remote
    edge above, and $u^{(0)}_k$ is the literal initial flux string.
    Thus the actual joint weight changes in the order
    \[
        \chi(p\eta_1^{-1}\eta_0)
        \longmapsto\chi(p\eta_1^{-1}k^{-1}\eta_0).
    \]
    The remote label is summed in the original-spin contraction.
    No tree gauge, factorization, Gram identity or state-action premise is
    supplied. This remains an auxiliary local crossing consequence:
    the identification of the complete prescribed geometric braid and
    its reunion experiment in \cite{Schuch2010PEPS}, lines 2560--2615,
    is separate.
\end{theorem}
\begin{proof}\leanok
    The remote bond tail has height $y$; every strip vertex has height
    $y+1$ or $y+2$. The period bound therefore puts that tail outside
    the actual region, including at a seam. Apply the derived uniform
    correlated column operation and the preceding actual cut identity.
    Extending the regional unitary by the complementary identity
    preserves unitarity.
\end{proof}
